%% Volet parallele : estimer la capacite d'un entrepot par imagerie
%% satellite, appele depuis la section 7.3.

\section{Capacity Observed Rather Than Declared}
\label{app:imagerie}

The whole measurement rests on the hypothesis stated in
Section~\ref{sec:realisation_experience}: that the weekly review is a sincere report, so the
indicators it carries and the urgency the same operator declares are two quantities rather
than one. The first campaign made it worth putting at least one of those indicators beyond
the declarant's reach altogether.

A second system was built alongside the main work to do exactly that: it
estimates the physical capacity of a warehouse from public satellite imagery, an observation
no operator supplies and none can bias. The two have not been connected, and no claim is made
that they have.

\subsection*{A physical limit, established by measurement}

The original objective was to count loading docks, and it cannot be done from a nadir view.
An empty dock door under a canopy produces no signal from directly overhead: measured contrast
at verified empty bays is $0.18\sigma$, $0.22\sigma$ and $0.01\sigma$ across a full sweep of
assumed depths, while occupied bays on the same facades register between $5.3\sigma$ and
$6.4\sigma$. The instrument is sensitive enough; the signal is absent. On 8~cm imagery the
doors stay equally invisible, which locates the cause in occlusion by the canopy roof rather
than in resolution: a gain of $2.5\times$ in ground sample distance recovers no additional
door.

Any method claiming to count empty doors one by one from overhead is therefore recovering an
artefact. Reaching that conclusion cost three weeks and saved the workstream from optimising a
quantity that is not in the image.

\subsection*{The reframing, and what actually moved the model}

The problem was restated as one the imagery can answer: delimit the equipped facade zone, and
derive the dock count from its length against a measured pitch. A segmentation model was
trained on hand-annotated imagery across four runs (Figure~\ref{fig:volume_training}).

\begin{figure}[htbp]
    \centering
    \includegraphics[width=0.86\linewidth]{Images/volume_training.pdf}
    \caption[Segmentation training: data helped, resolution did not]{Validation mask
    mAP50-95 across the four training runs. Quadrupling the annotated set roughly doubled the
    score; doubling the input resolution at a fixed data budget made every metric worse.}
    \label{fig:volume_training}
\end{figure}

Annotated volume drove the gain: quadrupling the training set from 86 to 351 images took mask
mAP50-95 from $0.155$ to $0.305$. Resolution did not. Run~4 is identical to run~3 except for a
doubled input size and is worse on every metric, because at a fixed data budget on a CPU-only
allowance the extra pixels cost more than they bought. The remaining work is therefore data
collection, not architecture search.

What is scored also had to change. Detection metrics measure shape overlap, but capacity
depends only on the length of the equipped zone, and a polygon drawn too deep still yields the
correct count. The project scores itself on the error in the derived count rather than on mAP.
On the reference run that gives a global bias of $-35\%$, a median error of $62\%$ on images
containing docks, and $89\%$ of true negatives correctly left empty: directionally useful, not
yet trustworthy, which is why it appears in Section~\ref{sec:perspectives} rather than in the
delivered scope.

\subsection*{An annotation supplier that reported its own work as verified}

An automated third-party annotation process returned 442 annotations and reported its own
output as verified and complete. It was unusable: it had placed a dock zone along every edge
of every building without examining the imagery, then subdivided each zone at the nominal
pitch, giving an implausible $88\%$ positive rate. All of it was discarded and the dataset
re-annotated by hand; an audit of the manual work found 11 problems across 687 zones, a defect
rate of $1.6\%$.

Checking cost one afternoon. Not checking would have cost a model trained to reproduce a
fiction, and every downstream number with it. The same discipline rejected a synthetic
logistics dataset in Section~\ref{sec:methode} and produced the baseline comparison of
Section~\ref{sec:temoin}.
